\subsection{作用}

定义1. 设 \(\Omega\) 和 \(\mathbf{E}\) 是两个集合。从 \(\Omega\) 到集合 \(\mathbf{E}^{\mathbf{E}}\) （ \(\mathbf{E}\) 到自身的映射所成的集合）的映射称为 \(\Omega\) 在 \(\mathbf{E}\) 上的一个作用.

设 \(\alpha \mapsto f_{\alpha}\) 是 \(\Omega\) 在 E 上的一个作用。映射 \((\alpha, x) \mapsto f_{\alpha}(x)\) （相应地 \((x, \alpha) \mapsto f_{\alpha}(x)\) ）称为与给定的 \(\Omega\) 在 E 上的作用相关联的 \(\Omega\) 在 E† 上的左（相应地，右）作用律。给定一个映射 \(g\) （从 \(\Omega \times E\) （相应地，E × \(\Omega\) ）到 E 中），存在且仅存在一个作用 \(\alpha \mapsto f_{\alpha}\) （ \(\Omega\) 在 E 上的），使得相关联的左（相应地，右）作用律为 \(g\) （集合论，II，§ 5，第 2 号，命题 3）。

在本章中，为简便起见，我们将说“作用律”

而不是“左作用律”。元素 \(f_{\alpha}(x)\) （E 的元素，对于 \(\alpha \in \Omega\) 和 \(x \in E\) ）有时称为 x 在 \(\alpha\) 下的变换，或 \(\alpha\) 与 x 的合成。它通常用左（相应地，右）乘法记号 \(\alpha.x\) （相应地 \(x.\alpha\) ）表示，点可以省略； \(\alpha\) 与 x 的合成此时称为 \(\alpha\) 与 x（相应地，x 与 \(\alpha\) ）的乘积。指数记号 \(x^{\alpha}\) 也被使用。在以下各段的论证中，我们通常使用记号 \(\alpha \perp x\) 。 \(\Omega\) 的元素常被称为算子。

例。(1) 设 E 是一个以乘法记法书写的结合广群。将严格正整数 n 与 E 到自身的映射 \(x \mapsto x^{n}\) 相关联的映射是 \(N^{*}\) 在 E 上的一个作用。如果 E 是一个群，那么将有理整数 a 与 E 到 E 的映射 \(x \mapsto x^{a}\) 相关联的映射是 Z 在 E 上的一个作用。

\begin{enumerate}
\def\labelenumi{(\arabic{enumi})}
\setcounter{enumi}{1}
\item
  设 \(\mathbf{E}\) 是一个律记为 \(\top\) 的广群。将 \(x \in \mathbf{E}\) 与映射 \(\mathbf{A} \mapsto x \top \mathbf{A}\) （ \(\mathbf{E}\) 的子集所成的集合到自身的映射）相关联的映射是 \(\mathbf{E}\) 在 \(\mathfrak{P}(\mathbf{E})\) 上的一个作用.
\item
  设 E 是一个集合。 \(E^{E}\) 的恒等映射是 \(E^{E}\) 在 E 上的一个作用，称为典范作用。相应的作用律是映射 \((f, x) \mapsto f(x)\) （从 \(E^{E} \times E\) 到 E 中）。
\item
  设 \((\Omega_{i})_{i\in I}\) 是一个集合族。对于所有 \(i\in I\) ，设 \(f_{i}\colon \Omega_{i}\to E^{\mathbb{E}}\) 是 \(\Omega_{i}\) 在 E 上的一个作用。设 \(\Omega\) 为诸 \(\Omega_{i}\) 的和（集合论，II，§4，第8号）。映射 \(f\) （从 \(\Omega\) 到 \(E^{\mathbb{E}}\) 上，扩张了诸 \(f_{i}\) ）是 \(\Omega\) 在 E 上的一个作用。这使我们能够将对一族作用的研究简化为对单个作用的研究。
\item
  给定 \(\Omega\) 在 \(\mathbf{E}\) 上的一个作用，其律记为 \(\bot\) ，一个子集 \(\Xi\) （ \(\Omega\) 的）和一个子集 \(\mathbf{X}\) （ \(\mathbf{E}\) 的）, \(\Xi \perp \mathbf{X}\) 表示 \(\alpha \perp x\) （其中 \(\alpha \in \Xi\) 和 \(x \in \mathbf{X}\) ）所成的集合；当 \(\Xi\) 由单个元素 \(\alpha\) 组成时，我们通常写 \(\alpha \perp \mathbf{X}\) 而不写 \(\{\alpha\} \perp \mathbf{X}\) 。将 \(\alpha \in \Omega\) 与映射 \(\mathbf{X} \mapsto \alpha \perp \mathbf{X}\) 相关联的映射是 \(\Omega\) 在 \(\mathfrak{P}(\mathbf{E})\) 上的一个作用，称为由给定作用通过扩张到子集集合而导出的作用。
\item
  设 \(\alpha \mapsto f_{\alpha}\) 是 \(\Omega\) 在 \(\mathbf{E}\) 上的一个作用。设 \(g\) 是 \(\Omega'\) 到 \(\Omega\) 的一个映射。则映射 \(\beta \mapsto f_{g(\beta)}\) 是 \(\Omega'\) 在 \(\mathbf{E}\) 上的一个作用.
\item
  设 \(f: \mathbf{E} \times \mathbf{E} \to \mathbf{E}\) 是集合 \(\mathbf{E}\) 上的一个合成律。映射 \(\gamma: x \mapsto \gamma_x\) （相应地 \(\delta: x \mapsto \delta_x\) ）（§ 2，第 2 号）将元素 \(x \in \mathbf{E}\) 与由 \(x\) 所作的左（相应地，右）平移相关联，它是 \(\mathbf{E}\) 在自身上的一个作用；它称为由给定律导出的 \(\mathbf{E}\) 在自身上的左（相应地，右）作用。当 \(f\) 是交换的时，这两个作用重合。
\end{enumerate}

与 \(\gamma\) 相关联的左（相应地，右）作用律是 f（相应地，f 的反合成律）。与 \(\delta\) 相关联的右（相应地，左）作用律是 f（相应地，f 的反合成律）。

设 \(\Omega, \mathbf{E}, \mathbf{F}\) 是集合， \(\alpha \mapsto f_{\alpha}\) 是 \(\Omega\) 在 \(\mathbf{E}\) 上的一个作用， \(\alpha \mapsto g_{\alpha}\) 是 \(\Omega\) 在 \(\mathbf{F}\) 上的一个作用。 \(\Omega\) -态射（ \(\mathbf{E}\) 到 \(\mathbf{F}\) 的），即 \(\mathbf{E}\) 到 \(\mathbf{F}\) 的与 \(\Omega\) 的作用相容的映射，是一个映射 \(h\) （从 \(\mathbf{E}\) 到 \(\mathbf{F}\) ），使得

\[
g _ {\alpha} (h (x)) = h \left(f _ {\alpha} (x)\right)
\]

对所有 \(\alpha \in \Omega\) 和 \(x \in E\) 成立。两个 \(\Omega\) -态射的复合是一个 \(\Omega\) -态射。

设 \(\Omega, \Xi, E, F\) 是集合， \(\alpha \mapsto f_{\alpha}\) 是 \(\Omega\) 在 \(E\) 上的一个作用， \(\beta \mapsto g_{\beta}\) 是 \(\Xi\) 在 \(F\) 上的一个作用， \(\phi\) 是 \(\Omega\) 到 \(\Omega'\) 的一个映射。一个 \(\phi\) -态射（ \(E\) 到 \(F\) 的）是一个映射 \(h\) （从 \(E\) 到 \(F\) ），使得

\[
g _ {\phi (\alpha)} (h (x)) = h (f _ {\alpha} (x))
\]

对所有 \(\alpha \in \Omega\) 和 \(x \in \mathbf{E}\) 成立 .

